% Options for packages loaded elsewhere
\PassOptionsToPackage{unicode}{hyperref}
\PassOptionsToPackage{hyphens}{url}
\PassOptionsToPackage{dvipsnames,svgnames,x11names}{xcolor}
\PassOptionsToPackage{space}{xeCJK}
\documentclass[
]{article}
\usepackage{xcolor}
\usepackage[margin=2.5cm]{geometry}
\usepackage{amsmath,amssymb}
\setcounter{secnumdepth}{-\maxdimen} % remove section numbering
\usepackage{iftex}
\ifPDFTeX
  \usepackage[T1]{fontenc}
  \usepackage[utf8]{inputenc}
  \usepackage{textcomp} % provide euro and other symbols
\else % if luatex or xetex
  \usepackage{unicode-math} % this also loads fontspec
  \defaultfontfeatures{Scale=MatchLowercase}
  \defaultfontfeatures[\rmfamily]{Ligatures=TeX,Scale=1}
\fi
\usepackage{lmodern}
\ifPDFTeX\else
  % xetex/luatex font selection
    \setmainfont[]{DejaVu Sans}
    \setmonofont[]{DejaVu Sans Mono}
  \ifXeTeX
    \usepackage{xeCJK}
    \setCJKmainfont[]{Noto Sans CJK SC}
          \fi
  \ifLuaTeX
    \usepackage[]{luatexja-fontspec}
    \setmainjfont[]{Noto Sans CJK SC}
  \fi
\fi
% Use upquote if available, for straight quotes in verbatim environments
\IfFileExists{upquote.sty}{\usepackage{upquote}}{}
\IfFileExists{microtype.sty}{% use microtype if available
  \usepackage[]{microtype}
  \UseMicrotypeSet[protrusion]{basicmath} % disable protrusion for tt fonts
}{}
\makeatletter
\@ifundefined{KOMAClassName}{% if non-KOMA class
  \IfFileExists{parskip.sty}{%
    \usepackage{parskip}
  }{% else
    \setlength{\parindent}{0pt}
    \setlength{\parskip}{6pt plus 2pt minus 1pt}}
}{% if KOMA class
  \KOMAoptions{parskip=half}}
\makeatother
\usepackage{longtable,booktabs,array}
\usepackage{calc} % for calculating minipage widths
% Correct order of tables after \paragraph or \subparagraph
\usepackage{etoolbox}
\makeatletter
\patchcmd\longtable{\par}{\if@noskipsec\mbox{}\fi\par}{}{}
\makeatother
% Allow footnotes in longtable head/foot
\IfFileExists{footnotehyper.sty}{\usepackage{footnotehyper}}{\usepackage{footnote}}
\makesavenoteenv{longtable}
\setlength{\emergencystretch}{3em} % prevent overfull lines
\providecommand{\tightlist}{%
  \setlength{\itemsep}{0pt}\setlength{\parskip}{0pt}}
\usepackage{bookmark}
\IfFileExists{xurl.sty}{\usepackage{xurl}}{} % add URL line breaks if available
\urlstyle{same}
\hypersetup{
  colorlinks=true,
  linkcolor={Maroon},
  filecolor={Maroon},
  citecolor={Blue},
  urlcolor={Blue},
  pdfcreator={LaTeX via pandoc}}

\author{}
\date{}

\begin{document}

\section{Generalization as Formalized Induction in the Unified Framework
of Representation, Logic, and Intuition --- Form Directs
Construction}\label{generalization-as-formalized-induction-in-the-unified-framework-of-representation-logic-and-intuition-form-directs-construction}

\section{统一框架 · 泛化作为归纳形式化 ·
形式指导构造}\label{ux7edfux4e00ux6846ux67b6-ux6cdbux5316ux4f5cux4e3aux5f52ux7eb3ux5f62ux5f0fux5316-ux5f62ux5f0fux6307ux5bfcux6784ux9020}

\begin{quote}
\textbf{Draft v0.1 (English)} \textbar{} 2026-08-11 Target:
NeurIPS/ICML/ICLR (theory + empirical) \textbf{Framework position}: Part
of the \emph{Unified Framework of Representation, Logic, and Intuition}
(manifesto:
\texttt{Unified\_Framework\_Representation\_Logic\_Intuition.md}). This
paper establishes the \textbf{form→construction} phase. Companion: P0
(three-channel equivalence), P2 (neural macro compilation). Evidence:
\texttt{docs/paper\_data/} (all experiments via run\_exp official judge
evaluation)
\end{quote}

\begin{center}\rule{0.5\linewidth}{0.5pt}\end{center}

\textbf{Repository:
\href{https://github.com/YuchenWang-ai/Generalization-as-Formalized-Induction-Token-Native-Factorized-Representation-and-Weight-Compilation}{YuchenWang-ai/Generalization-as-Formalized-Induction-Token-Native-Factorized-Representation-and-Weight-Compilation}}

\subsection{Abstract}\label{abstract}

We propose that systematic generalization is not an emergent or a priori
faculty of neural networks, but the formalization of induction. To
generalize on a concept A, A must be precisely defined, and every
concept in A's definition chain must be sufficiently trained; sufficient
training is a relational property --- simultaneous training of parallel
concepts under the same superordinate concept, wrong-answer examples,
and collocation coverage --- not a quantitative one. We operationalize
this in a token-native representation (six projection layers: axioms B,
concepts C, symbols S, grammar G, presentation P, arrows A) in which
conceptual degrees of freedom are factorized according to their
symmetries, and execution is driven by token definitions (a
zero-hardcoding evaluator). A 2-layer, \textless1 MB transformer
generalizes to arbitrary width (2000-digit addition, acc 1.000), unseen
radices ({[}3..9{]}, OOD 0.998), unseen domains, and joint Cartesian OOD
--- with \textasciitilde2000 samples/3 epochs for a simple arithmetic
suite and \textasciitilde6000 samples/8 epochs for the full
logic-and-arithmetic suite. Symbol permutation invariance holds
(0.996→0.987 under 26080 digit renamings), and the three execution
channels --- construction, intuition, formal rewriting --- are
semantically equivalent (1.000 on 2000-digit carry chains), constituting
the neural implementation of Computational Trinitarianism. We emphasize
that intuition is the reasoning fast path: it is hard to obtain
(acquisition requires config-dependent epochs --- 3 for a simple suite,
8 for the full logic-and-arithmetic suite), but once formed it
generalizes easily by skipping steps, jumping from input to output over
the intermediate construction steps, with zero decay out to 2000-digit
width, radix 60, and 100\% digit anonymization. Acquisition and
step-skipping generalization are decoupled. Per-concept learning curves
are sharply bimodal --- a concept is either compiled (100\% by epoch 3)
or not (0\%, unfixed by more epochs) --- the mechanism-level signature
that definition-precise grammar yields extremely fast training
(form→construction), and that definitionally-impoverished tokens never
generalize. Within the Unified Framework, this paper is the
\textbf{form→construction} leg: correct syntax yields extreme training
speed; correct construction yields extreme generalization (OOD); and
correct generalization compiles into extreme intuition (the fast
reasoning path).

\subsection{1 Introduction}\label{introduction}

\textbf{The problem.} Compositional generalization (Lake et al., 2017;
Keysers et al., 2020) is a well-documented phenomenon: models sometimes
respond correctly to combinations absent from training. Yet the
literature asks \emph{whether} models generalize, not \emph{how}
generalization arises. Three defaults dominate: (i) generalization is an
emergent/a priori faculty; (ii) generalization follows from data
engineering at scale; (iii) the relevant structure must be provided by
an external program (neural-symbolic hybrids, program synthesis). Each
default precommits to a mechanism without explaining the observed
pattern. In the philosophy of mathematics this is the old dispute
resurfaced: Platonism (structure exists a priori, to be discovered),
intuitionism (existence = constructibility, Brouwer), and formalism
(mathematics = sign games, Hilbert) --- each precommits a source for
structure.

Meanwhile, the dominant approach --- ``existing mathematical language,
then tokenize, then train'' --- pre-supposes primitives. λ-calculus
fixes variable/abstraction/application (Barendregt 1984); Church
encoding \emph{implements} any concept but erases all structural
boundaries in the representation --- a model trained on the encoded
surface can memorize the function table but has no access to the
structure that would make out-of-distribution combinations easy. A
concrete instance is the ``neural λ-calculus'' program (Flach, Moreira
\& Lamb 2023): a seq2seq Transformer is trained on standard λ-syntax
with de Bruijn indices to perform β-reduction, reaching 97.7\% on
Boolean terms --- but generalization is confined to the λ-term
distribution it was trained on; there is no width/radix/order
extrapolation, no symbol permutation, no concept formation. This is
precisely the ``λ-tokenize-then-train'' route this paper argues is
incomplete. Our internal structural analysis of λ-iteration
(f\^{}\{m+n\} = f\^{}m ∘ f\^{}n; the Church zero is not a base;
iteration becomes a universal property under initiality) shows that the
relevant structure is genuinely present in mathematics, but is
\emph{erased} by the encoding --- it must be re-exposed in the
representation for a learner to use it.

\textbf{What we do.} We take the opposite route. We treat systematic
generalization as an engineering target: it is produced, not assumed.
The producing mechanism is the formalization of induction --- a concept
A generalizes exactly when (i) A is precisely defined and (ii) every
concept in A's definition chain is sufficiently trained, where
sufficiency is relational (parallel concepts, wrong answers,
collocations). We build a token-native representation in which these
conditions are verifiable, and show that a minimal transformer then
generalizes essentially for free.

\textbf{Contributions} (each a falsifiable claim; the propositions G0,
G1, G2, G4, G5 and R are stated in §3):

\begin{itemize}
\tightlist
\item
  \textbf{C1 (Theory) --- Generalization as formalized induction.}
  Systematic generalization is not emergent/a priori: definition
  completeness (G1) and relational sufficiency (G2) are sufficient. The
  anti-prior claim is empirically falsifiable
  (definitionally-impoverished tokens do not generalize even with
  abundant data).
\item
  \textbf{C2 (System) --- Token-native factorized representation.} Six
  projection layers (B/C/S/G/P/A) with degrees of freedom factorized by
  symmetry, and a zero-hardcoding evaluator; ``definition precision''
  becomes checkable (no dangling references, cycle certificates,
  coverage audits).
\item
  \textbf{C3 (Empirics) --- Minimal-model systematic generalization.} A
  2-layer, \textless1 MB transformer: arbitrary width (2000-digit,
  1.000), unseen radices (OOD 0.998), unseen domains, joint OOD;
  symbol-permutation invariance (0.987); three-channel semantic
  equivalence (1.000). Sample counts: \textasciitilde2000 (simple) /
  \textasciitilde6000 (full suite); epochs 3 (simple) / 8 (full).
\item
  \textbf{C4 (Mechanism) --- Weight compilation.} Training compiles the
  definition-driven construction channel into weights; supervision
  excludes answer traces, so the model learns structure rather than
  replicating the unfolding. The product is \emph{intuition --- the
  reasoning fast path}: hard to obtain, easy to generalize by
  step-skipping.
\item
  \textbf{C5 (Epistemology) --- Three-channel equivalence.} Construction
  / intuition / formal rewriting are pairwise semantically equivalent on
  the same inputs --- the neural implementation of Computational
  Trinitarianism (equivalence between execution channels of one system,
  not between logics).
\end{itemize}

\subsection{2 Problem Setting}\label{problem-setting}

\textbf{Systematic generalization.} The ability to respond correctly to
combinations not present in training, where combinations span
independent degrees of freedom: width × radix × operator order × domain
× nesting depth. We test extrapolation in these dimensions, never
interpolation.

\textbf{In-distribution vs OOD.} A degree-of-freedom value is
in-distribution if its combination appears in training; OOD otherwise.
Single-dimension extrapolation (width beyond training length; radix
absent from training; order beyond the highest trained) and joint
extrapolation (unseen points of the Cartesian product) are both tested.
Primary metric: \emph{judgment accuracy} --- full output-sequence
reconstruction of the final truth token (the only trusted metric;
positional accuracy can mask structural errors, e.g.~a 0.904 positional
score with 0.000 judgment accuracy on a focused-configuration artifact).
All evaluation uses the single official run\_exp judge path.

\textbf{Notation and conventions.} A \emph{concept} (token) is the unit
of meaning; a \emph{definition} is a structured statement of a concept's
content, among five forms (explicit, inductive, recursive, axiomatic,
implicit). A \emph{definition chain} of A is the transitive closure of
concepts that A's definition references. \emph{Sufficient training} of a
concept means relational coverage: it participates with its parallel
concepts, its negative examples, and its collocations. \emph{Judgment}
is a complete output sequence (expression, result, truth);
\emph{judgment accuracy} is the fraction of full reconstructions
correct.

\textbf{Grammar.} Expression grammar is data-driven (G-layer gtoken
arrangement + P-layer presentation), with slots arg/fn/args/binder/body
and node types
atom/application/equality/binary\_connective/unary\_connective/quantified/eos.
Numerals are a sign/radix/cardinality intrinsic vector combined with a
digit sequence; any radix is representable (no single-character names
needed). Propositions are judged along gtoken/ptoken-assembled
sequences.

\subsection{3 Theory: Formalized
Induction}\label{theory-formalized-induction}

We state the theory as six falsifiable propositions (G0, G1, G2, G4, G5
and R). Each is a claim about when and how generalization is produced;
each is supported by experiments in §6 and ablations in §8.

\textbf{G0 (epistemic ground) --- Generalization is the formalization of
induction.} \emph{There is no a priori or transcendental generation of
conceptual ability: any ability to generalize on a concept must be
produced, not assumed.} Empirically falsifiable: if concept ability were
a priori, definitionally-impoverished concepts would still generalize;
they do not --- an impoverished vocabulary requires an order of
magnitude more training to converge (40 vs 3 epochs, §8), and genuinely
under-defined tokens never generalize. EXP-61 quantifies the boundary:
in the full suite, removing definition-form samples (law+definition)
leaves e3 accuracy nearly unchanged (exact 0.295 vs blur 0.271) and both
variants converge by e10 --- the definition-precision leverage on
training speed is decisive only under extreme deficiency
(whole-vocabulary poverty), not under partial removal of definition
samples. Grounding: this is the anti-Platonist / anti-nativist position
(Brouwer's constructivism in machine form; Fodor 1975; Chomsky 1959 as
the opposing positions).

\textbf{G1 (definition completeness).} \emph{To generalize on concept A,
(i) A's definition must be precise and complete, and (ii) every concept
in A's definition chain must itself be sufficiently trained. Both
conditions are necessary; neither suffices alone.} Evidence: a
definition-chain token with a single training instance (value\_three) is
not learned until coverage is added; a definitionally-impoverished
vocabulary requires 40 epochs vs 3 for a precise one (§8). Definition
forms follow Aczel 1977, Gupta 2023, Belnap 1993.

\textbf{G2 (relational sufficiency).} \emph{Sufficient training is a
relational property, not a quantitative one: (i) simultaneous training
of parallel concepts under the same superordinate concept makes
relations between concepts explicit learning signals; (ii) wrong-answer
training (negative examples, including structure-identical-result-wrong)
defines what a concept is not; (iii) collocation training covers each
participating token's adjacent collocations above a threshold.}
Evidence: removing one gate-family's supervision collapses that gate to
0.000 and \emph{degrades all other gates from 1.0 to 0.833} (§8);
increasing negative-example coverage lifts a relation from 0.000 to
0.958.

\textbf{G4 (modal separation).} \emph{Mirroring
perception--meaning--expression, model input, logic, and output are
assigned independent grammars; signifier and signified are separated at
the token level.} Evidence: assigning one symbol to two semantic
wrappers collapses training (§8). Grounded in semiotics (Saussure 1916;
Frege 1892) and in abstract/concrete syntax separation.

\textbf{G5 (intuition/reasoning separation).} \emph{Intuition is the
reasoning fast path: hard to obtain (acquisition requires
config-dependent epochs), but easy to generalize by skipping steps ---
it jumps from input to output over the intermediate construction steps.}
Evidence: acquisition needs 8 epochs for the full suite (exp53) while
step-skipping generalization shows zero decay at 2000 digits, radix 60,
and 100\% anonymization (exp80) --- acquisition and step-skipping
generalization are decoupled (§6).

\textbf{G5b (positional sensitivity separates intuition from logic;
2026-08-11).} \emph{The intuition channel is position-insensitive
(robust to token/positional scrambling); the logic channel is
position-sensitive (requires a determinate position-to-semantics
mapping).} This follows from two observations. First, EXP-80's 100\%
digit anonymization at 1.000 shows the compiled fast path ignores
symbol/position detail and extracts structural relations. Second,
EXP-61h (intervention probing, §8) shows that the logic channel
\emph{cannot} tolerate position ambiguity: with two synonym tokens at
one operator position the model canonicalizes to a single token (V2same:
c0=1.000, c1=0.000) --- structure demands a determinate position; with
three candidates at one position (V3diff/V3mix) judgment collapses
(0.250). Position scanning confirms the intuition side: under V3mix (two
synonym and tokens + one or token), the surviving or token scores 0.750
at every one of three operator positions ({[}op{]}AB, A{[}op{]}B,
AB{[}op{]}) --- the retained semantics is position-insensitive. This
explains a familiar cognitive phenomenon: scrambling the letters of a
sentence does not impede understanding (the intuition channel is
position-insensitive), yet \emph{deliberately attending} to the
scrambled text to re-derive its logic makes it lose meaning (the logic
channel requires determinate positions). Intuition is compiled structure
(position-free); logic is unfolding structure (position-bound) --- a
mechanistic reading of G5.

\textbf{R (three-channel equivalence).} \emph{A single system exhibits
three execution channels for the same mathematical object --- the
construction channel (unfolded, definition-driven expansion), the
intuition channel (compiled weight execution), and the formal channel
(definition-driven rewriting) --- and they are pairwise semantically
equivalent: construction can be compiled into intuition, and both agree
with formal rewriting.} This is the neural implementation of
Computational Trinitarianism (Curry--Howard: programs = proofs = types),
with equivalence holding between execution channels of one system rather
than between logics. Evidence: at a 2000-digit carry chain, construction
= formal exactly and intuition scores 1.000 (§6).

\subsection{4 The Representation
System}\label{the-representation-system}

\textbf{Six projection layers.} A token is one object; the six JSON
projections are projections of it.

\begin{longtable}[]{@{}
  >{\raggedright\arraybackslash}p{(\linewidth - 4\tabcolsep) * \real{0.3333}}
  >{\raggedright\arraybackslash}p{(\linewidth - 4\tabcolsep) * \real{0.3333}}
  >{\raggedright\arraybackslash}p{(\linewidth - 4\tabcolsep) * \real{0.3333}}@{}}
\toprule\noalign{}
\begin{minipage}[b]{\linewidth}\raggedright
Layer
\end{minipage} & \begin{minipage}[b]{\linewidth}\raggedright
Entity
\end{minipage} & \begin{minipage}[b]{\linewidth}\raggedright
Role
\end{minipage} \\
\midrule\noalign{}
\endhead
\bottomrule\noalign{}
\endlastfoot
B axioms & baseloop & axiomatic endpoints; terminal of definition-chain
tracing \\
C concepts & derive & five definition forms
(explicit/inductive/recursive/axiomatic/implicit) \\
S symbols & symbol & signifiers (glyph → maps\_to many-to-many;
ambiguity preserved, resolved by context) \\
G grammar & gtoken & arrangement methods (slots/reduction); syntax
itself is tokens \\
P presentation & presentation & concrete notation
(symbol/precedence/associativity); one concept, many notations \\
A arrows & arrow & directed relations between concepts (domain lifts,
iteration chains, duals) \\
\end{longtable}

Transformations depend only on tokens; knowledge lives in token/grammar
data; code reads and executes. This is the zero-hardcoding discipline.
The six-layer projection is a formalization of representation rather
than a new logic: it separates signifier from signified (Saussure 1916;
Frege 1892), makes arrangement a definable object (van Wijngaarden 1965;
Ford 2004; Rendel \& Ostermann 2010), separates binding structure from
syntax (de Bruijn 1972; Fiore--Plotkin--Turi 1999; Gabbay \& Pitts
2002), and encodes concept relations as first-class arrows (Goguen et
al.~1977; Mac Lane 1971).

\textbf{Degrees of freedom.} A concept is a coordinate (d1,\ldots,dk),
not an isolated token: iteration order (succ→+→×→\^{}→↑↑),
duality/direction (inverse: succ/pred, power/root, +/−), domain
(natural/integer/rational/real/complex), and notation
(sign/base/cardinality).

\textbf{Symmetry.} Iterative symmetry: iterate(direction, level,
operator) is a higher-order operator; F\_\{n+1\}(a,b) = Iterate(F\_n, a,
b). Structurally, iteration is composition (f\^{}\{m+n\} = f\^{}m ∘
f\^{}n): it yields a monoid, then a category, then a functor and natural
transformation under semiconjugacy --- the standard route from
λ-iteration to category theory (Lambek 1968; Mac Lane 1971; our λ-base
iteration study). Dual symmetry: inverse arrows (power↔root), symmetry
families (neg/scale/root; complement/parallel\_sum as De Morgan duals;
translation/inversion as modular group). Domain lifts: coercion as
arrows (retrieved by concept=coercion, zero hardcoding), in the style of
canonical inclusions between algebraic structures (Mac Lane 1971; Adámek
et al.~2004).

\textbf{Composition and iteration.} Arrangement lives in the G-layer;
assembly reads data only; SKI combinators are gtokens (reduction rules
in data). Iteration is an operator over operators; iteration-chain
membership is discovered along A-layer arrows (ARROW\_REGISTRY), never
via name hardcoding. The generic evaluator handles pure chain members by
iterative reduction and special-symmetry members
(translation/inversion/differential) by delegated symmetric evaluation
--- token in, token out. This mirrors the definitional structure of
iteration (operational → universal property under initiality; Lambek
1968; Goguen et al.~1977), kept as data rather than code.

\textbf{Domain/radix representation.} The radix is an \emph{object}, not
a structural constant: a numeral is {[}sign, base, cardinality{]} ×
digit vector; one set of weights handles any base (base as input
parameter). This follows the standard analysis of positional notation
(base as a parameter of representation, not of the number --- see our
number-system survey); notations are extensible
(fraction/radical/imaginary\_unit) while numeral token sequences remain
digit-arranged (never ``directly by value''). Different semantic
wrappers require different symbols --- operation nesting ``(\ldots)'' vs
representation nesting ``{[}\ldots{]}'' --- a modal separation whose
violation collapses training (see §8); this echoes the parse/render
separation between concrete and abstract syntax (Rendel \& Ostermann
2010; Fiore--Plotkin--Turi 1999).

\subsection{5 Weight Compilation}\label{weight-compilation}

\textbf{Data generation.} Samples are instantiations of definitions.
Concepts carry one of five definition forms (explicit, inductive,
recursive, axiomatic, implicit), following the standard theory of
definitions (Aczel 1977; Gupta 2023; Belnap 1993; Suppes 1957):
inductive definitions rest on least-closure, implicit definitions on
conservativeness and eliminability. The synthesizer (config-driven)
composes: balanced / trace / numeral\_split / fill / choose / definition
/ arrow / law / logic\_nested / logic\_arith / logic\_interdef /
coercion / deep\_nest / extrap / cartesian / radix / neg. Sufficiency
(G2) requires: parallel concepts trained together (relations become
learning signals), wrong answers (negative examples, including
structure-identical-result-wrong), and collocation coverage
(adjacent-pair coverage ≥ 3). Sequences are never hand-written; assembly
follows gtoken/ptoken (hand-written sequences are flagged as
``missing-token'' signals).

\textbf{Reference evaluator.} The generic evaluator computes from token
definitions: logic/compare truth tables, arithmetic iteration chains,
symmetry families, primality --- all read from definitions. Construction
channel (token\_iterator) and formal channel (rewriting engine) share
the same definition source, which is the basis of three-channel
consistency.

\textbf{Training.} A 2-layer, dim-64 transformer (\textless1 MB). Two
objectives: positional cross-entropy on positive samples, and a validity
cross-entropy that separates correct from incorrect judgments (the model
must not merely copy the input --- it must judge). Sample counts and
convergence are config-dependent: \textasciitilde2000 samples / 3 epochs
for the simple arithmetic suite; \textasciitilde6000 samples / 8 epochs
for the full logic-and-arithmetic suite (exp53 curve).

\textbf{Evaluation discipline.} Judgment accuracy --- full
output-sequence reconstruction of the final truth token --- is the only
trusted metric; positional accuracy can mask structural errors. A single
official evaluation path is used for every experiment (including
ablations), because a divergence in evaluation branches silently
invalidates comparisons. Position coverage is treated as an independent
degree of freedom (a syntactic slot never seen in training does not
generalize even with complete definitions). Diagnostics: per-token
generalization (bimodal, §6), epoch curves, crash-culprit (which token
most often breaks a judgment), and coverage audits. These make
``definition completeness'' and ``relational sufficiency'' checkable at
the token level rather than as aggregate scores.

\textbf{Why inputs exclude answer traces.} Supervision contains only
judgment truth (correct/incorrect), never the step-by-step unfolding. If
traces were given, the model could copy the unfolding (sequence
cloning); without them it must induce the computational structure from
definitions and compile it into weights. Learning is compilation. The
compiled product is intuition: the reasoning fast path.

\subsection{6 Experiments}\label{experiments}

All runs use the single official run\_exp pipeline
(config/metrics/model/samples/views archived). Judgment accuracy
reported.

\begin{longtable}[]{@{}
  >{\raggedright\arraybackslash}p{(\linewidth - 2\tabcolsep) * \real{0.5000}}
  >{\raggedright\arraybackslash}p{(\linewidth - 2\tabcolsep) * \real{0.5000}}@{}}
\toprule\noalign{}
\begin{minipage}[b]{\linewidth}\raggedright
Experiment
\end{minipage} & \begin{minipage}[b]{\linewidth}\raggedright
Result (judgment acc)
\end{minipage} \\
\midrule\noalign{}
\endhead
\bottomrule\noalign{}
\endlastfoot
Full suite, seed s1/s2 (9 logic gates + arithmetic) & 1.000 / 1.000
(\textasciitilde6000 samples, 8 epochs) \\
Simple arithmetic suite (arith\_multi\_type) & 1.000
(\textasciitilde2000 samples, 3 epochs) \\
imply with supervision (definition+position+judgment) & 0.996 \\
2000-digit addition extrapolation & 1.000 \\
Unseen radix {[}3..9{]} × 20 digits & OOD 0.998 \\
Unseen domain / operator level (addition↔multiplication) & generalizes;
hyper-operator order coverage incomplete (see Limitations) \\
Joint Cartesian OOD (width × radix × nesting) & 1.000 (exp80 joint) \\
Symbol permutation (exp41: 26080 digit renamings) & 0.996→0.987 \\
Three-channel consistency (exp50: 2000-digit carry chain) & 1.000 \\
Fast/slow path cost (exp51: 2000 digits) & construction 0.62s /
intuition 0.086s (7×) / formal 0.069s (9×) \\
Convergence curve (exp53: epochs 1/2/3/5/8) &
0.001/0.214/0.244/0.781/0.996 \\
Step-skipping zero-decay (exp80: 2000-digit / radix 60 / 100\%
anonymized / joint) & 1.000 \\
Step-skipping machine form (exp90: 20-digit add) & construction 82
tokens/0.410ms → intuition 5 tokens/0.186ms (2.2×; honest lower bound,
itoken outside vocab) \\
State propagation (exp70: 2000-digit carry and borrow chains) & 1.000
--- real per-digit state propagation, weights ≈ program \\
Position sensitivity (exp61h: synonym tokens at one op position) & 2
tokens → canonicalize (c0 1.000 / c1 0.000); 3 tokens → collapse (0.250)
--- logic is position-sensitive \\
Arrangement freedom (notation: prefix vs infix vs mixed) & uniform
notation → identical learnability (1.0 full; 0.677 low-training);
dual-position mixing (same op prefix+infix) improves training (0.948 vs
0.852); imply judge-position coverage +0.019 --- arrangement is free,
coverage decides \\
\end{longtable}

\textbf{Key experiments.}

\begin{itemize}
\tightlist
\item
  \emph{Width}: training width ≤ L, test 2000 digits ---
  length-independence, refutes function-table memorization.
\item
  \emph{Radix}: base is an input parameter, not a structural constant
  --- refutes decimal memorization.
\item
  \emph{Order/domain}: addition↔multiplication level extrapolation
  succeeds; hyper-operator levels (power/root/tetration) are not fully
  covered because radix-base generation overflows for astronomical
  values --- see Limitations. Iteration-chain structure is learned via
  A-layer arrows, not per-operator tables.
\item
  \emph{Symbol permutation (exp41)}: random permutation of 10 digit
  concept tokens (26080 renamings, structure preserved), tested on
  original tokens --- invariance 0.996→0.987. The model learned
  structural relations (successor/iteration/carry), not digit semantics.
  Combined with exp80's 100\% anonymization at 1.000, this is the
  strongest refutation of ``mathematical knowledge as memory''.
\item
  \emph{Three channels (exp50)}: construction = formal exactly
  (2/9/20/2000 digits); intuition 1.000 including a 2000-digit carry
  chain (999\ldots9+1) --- real state propagation, weights ≈ program.
\item
  \emph{Step-skipping zero decay (exp80)}: at fixed acquisition (8
  epochs), extrapolation distance increases with no decay --- 2000-digit
  (10× width), radix 60 (6×), 100\% digit anonymization, and joint
  width×radix all at 1.000. Acquisition and step-skipping generalization
  are decoupled.
\item
  \emph{Step-skipping machine form (exp90)}: the intuition path treats a
  20-digit number as an atomic num token (5 tokens) versus the
  construction path's 82-token digit-wise unfolding. Honest caveat:
  itoken is outside the model vocab (pad id 0); the measurement is the
  sequence-length effect on forward cost (a lower bound). If the model
  were to truly learn itoken semantics, the intuition path would be an
  O(1) decision with a larger advantage.
\end{itemize}

\textbf{Per-token generalization: extreme bimodal distribution.} The
aggregate learning curve (exp53) is the union of per-token curves that
are sharply bimodal:

\begin{itemize}
\tightlist
\item
  A learned concept reaches 100\% by epoch 3 with no gradual transition;
\item
  an unlearned concept stays at 0\% and is \emph{not} fixed by more
  epochs;
\item
  intermediate states (30\%, 50\%) are essentially absent.
\end{itemize}

The apparent gradual rise of the aggregate curve is therefore the
stacking of per-token ``switch-ons'': each concept is either compiled
into weights (100\%) or not (0\%). This is the mechanism-level signature
of the compilation account (C4). It also makes generalization failures
attributable: an unlearned item points at the concept itself (its
definition precision or its dependency/coverage), not at ``insufficient
training'' --- consistent with C1 and with the token-level diagnostics
(per-token accuracy, crash-culprit, coverage audit) used throughout.

\subsection{7 Semantic-Path
Compilation}\label{semantic-path-compilation}

\textbf{Zero samples already execute.} A new operator σ's semantic
closure is guaranteed by its definition chain: a base model that has
never seen σ can still execute along the expansion path T\_σ
(construction channel) --- correct, but long. Semantic closure does not
depend on training.

\textbf{Few samples build the reasoning fast path.} A small number of
samples compiles a semantically equivalent intuition channel:
Sem\_fast(σ) = Sem\_expanded(T\_σ) (verifiable, with fallback);
Cost\_fast \textless{} Cost\_expanded. The intuition channel does not
expand the semantic closure --- it solves an execution-cost problem, not
a new-semantics problem. Unlike classical distillation, the teacher and
student are the \emph{same model's} construction and intuition channels
(construction-channel self-compilation), with a verifier and fallback.

\textbf{Terminology.} Intuition is the \emph{reasoning fast path}. It is
hard to obtain (sufficient training; config-dependent epochs), but easy
to generalize by skipping steps: it jumps from input to output over the
intermediate construction steps. Speed is its nature as a reasoning
path, not an independent optimization; step-skipping is the mechanism of
its outward generalization. Honest caveat: the formal engine (eval\_op)
is faster in absolute terms (0.069s vs 0.086s at 2000 digits); the value
of the intuition channel is end-to-end execution (no external oracle)
and preserved step-skipping generalization.

\subsection{8 Ablations}\label{ablations}

\begin{longtable}[]{@{}
  >{\raggedright\arraybackslash}p{(\linewidth - 4\tabcolsep) * \real{0.3333}}
  >{\raggedright\arraybackslash}p{(\linewidth - 4\tabcolsep) * \real{0.3333}}
  >{\raggedright\arraybackslash}p{(\linewidth - 4\tabcolsep) * \real{0.3333}}@{}}
\toprule\noalign{}
\begin{minipage}[b]{\linewidth}\raggedright
Ablation
\end{minipage} & \begin{minipage}[b]{\linewidth}\raggedright
Result
\end{minipage} & \begin{minipage}[b]{\linewidth}\raggedright
Supports
\end{minipage} \\
\midrule\noalign{}
\endhead
\bottomrule\noalign{}
\endlastfoot
Degrees of freedom: factored (numeral\_v4 correct) vs
entangled/wrong-symmetry & entangling collapses generalization &
G4/C1 \\
Parallel-concept removal (exp20a/b: delete iff/xor or and/or
supervision) & overall 0.798/0.814; deleted gate 0.000; \textbf{other
gates 1.0→0.833} --- the mutual-training chain (De Morgan duals /
inverse arrows) breaks & G2 \\
Symbol permutation (exp41) & 0.996→0.987 & G4 \\
Three-channel consistency (exp50) & construction=formal=intuition
(2000-digit carry chain, 1.000) & R/C5 \\
Convergence curve (exp53) & monotone; full suite converges at 8 epochs
--- supports ``intuition is hard to obtain'' & G5 \\
Negative-example coverage (neg 9→15) & neg 0.000→0.958 & G2 \\
Definition-chain token uncovered (value\_three, 1 sample) & not learned;
learned after coverage & G1 \\
Definitionally impoverished (vocab 153) & requires 40 epochs & G0/G1 \\
Definition-form ablation (EXP-61: exact vs blur, remove law+definition)
& e3 0.295 vs 0.271 (exact marginally higher); \textbf{both e10→1.000}
--- in the full suite, definition structure gives marginal (not
decisive) training-speed gain; the decisive gain requires extreme
deficiency (vocab 153) & G0/G1 \\
Wrapper-symbol mixing (``( )'' vs ``{[} {]}'' same symbol) & training
collapses & G4 \\
Judgment vs positional (focused config) & positional 0.904 → judgment
0.000 (artifact) & metric discipline \\
Position ambiguity (EXP-61h: 2 synonym tokens at one op position) &
canonicalizes to single token (c0=1.000, c1=0.000); 3 candidates →
collapse (0.250) --- logic is position-sensitive & G5b \\
\textbf{Ordinary representation baseline} (EXP-C1: radix-fixed /
iteration-hidden / entangled) & pending & C1 \\
\end{longtable}

\textbf{The ordinary-representation baseline (EXP-C1).} The decisive
control: same model, same data budget, but the representation is
de-factored. We progressively destroy degree-of-freedom factorization
and observe extrapolation collapse:

\begin{longtable}[]{@{}
  >{\raggedright\arraybackslash}p{(\linewidth - 4\tabcolsep) * \real{0.3333}}
  >{\raggedright\arraybackslash}p{(\linewidth - 4\tabcolsep) * \real{0.3333}}
  >{\raggedright\arraybackslash}p{(\linewidth - 4\tabcolsep) * \real{0.3333}}@{}}
\toprule\noalign{}
\begin{minipage}[b]{\linewidth}\raggedright
Control
\end{minipage} & \begin{minipage}[b]{\linewidth}\raggedright
Factorization destroyed
\end{minipage} & \begin{minipage}[b]{\linewidth}\raggedright
Result
\end{minipage} \\
\midrule\noalign{}
\endhead
\bottomrule\noalign{}
\endlastfoot
C1d factored (exp02, reference) & none & zero-decay extrapolation
(exp80: 1.000) \\
C1b iteration-hidden (done) & iterate structure not exposed (law
power/root/tetration, iterate/inverse arrows, staircase removed) &
overall 0.745 (Δ0.25); \textbf{root 0.000, tetration 0.000} ---
high-level operators collapse; \textbf{power 1.000} (derivable from
multiplication: 2\^{}3 = 2×2×2); addition/multiplication 1.000 \\
C1a radix-fixed (re-designing) & radix made a structural constant (train
base-10 only, no base/cardinality parameter) & pending --- expected:
unseen radix (base 16/60) collapses; first attempt (strip base tokens)
over-destroyed the numeral representation (0.003), being re-designed to
keep base-10 computable while removing the radix degree of freedom \\
C1c entangled (pending) & composite concepts merged into single tokens
(degrees of freedom fused) & pending --- expected: new combinations
collapse \\
\end{longtable}

\textbf{C1b analysis.} The iteration-hidden result is precise evidence
for the definition-chain account (G1): collapsing operators are exactly
those whose definition chains genuinely depend on the iteration
structure (root = inverse of power-iteration; tetration =
power-iteration at the next level), while power survives because it is
reducible from multiplication. Factorization of the iteration degree of
freedom is thus the \emph{cause} of order extrapolation --- not for all
operators uniformly, but exactly for those whose definitions require it.
This is a sharper claim than ``factorization helps generally.''

\textbf{Methodological discipline (from a withdrawn experiment).} An
earlier zero-supervision implication experiment (EXP-10, 0-IMPLY) was
fully withdrawn: its reported ``semantic emergence'' (bare truth table
0.75 / deep nesting 0.938) was an artifact of last-token evaluation plus
syntactic misplacement (an evaluation-branch inconsistency that made all
variants produce identical results). Under the official full-sequence
judge, implication scored 0.000; root cause: positional misplacement
plus insufficient sample count. We record the lessons --- full-sequence
reconstruction is the only trusted metric; evaluation must use the
single official pipeline; position coverage is an independent degree of
freedom --- and the correct path for implication is the triple
definition+position+judgment supervision (0.996). The withdrawn
experiment contributes methodological discipline, not evidence.

\subsection{9 Related Work}\label{related-work}

\begin{enumerate}
\def\labelenumi{\arabic{enumi}.}
\tightlist
\item
  \textbf{λ-calculus and formal-system families.} Interpreted object,
  not benchmark: primitives pre-supposed
  (variable/abstraction/application); Church encoding erases structure.
  The ``neural λ-calculus'' program (Flach, Moreira \& Lamb 2023)
  exemplifies the λ-tokenize-then-train route --- a seq2seq Transformer
  trained on de Bruijn-indexed λ-syntax to perform β-reduction, with
  generalization confined to the training distribution. The structural
  reading of iteration --- that iteration is composition (f\^{}\{m+n\} =
  f\^{}m ∘ f\^{}n), that the Church zero is not a base, and that
  category theory turns iteration into a universal property --- is
  developed in detail in our internal study (λ-base iteration → category
  theory); standard foundations: Lambek 1968 (fixpoint theorem for
  complete categories), Goguen et al.~1977 (initial algebra semantics),
  Martin-Löf 1984 (intuitionistic type theory), Fiore--Plotkin--Turi
  1999 (abstract syntax and variable binding).
\item
  \textbf{Definition theory and definition forms.} Five definition forms
  (explicit / inductive / recursive / axiomatic / implicit) follow Gupta
  2023 (SEP ``Definitions'') and Belnap 1993 (``On Rigorous
  Definitions''); inductive definitions per Aczel 1977; implicit
  definitions require conservativeness and eliminability (Suppes 1957;
  Gupta--Belnap 1993 revision theory); definitional machinery in proof
  assistants (Lean, Rocq, Isabelle). No single prior system unifies
  concept definitions and grammar definitions in one schema --- this is
  a gap this work fills.
\item
  \textbf{Semiotics and signifier/signified separation.} Saussure 1916
  (arbitrariness, differential principle); Peirce's sign triad; Frege
  1892 (Sinn/Bedeutung); Harnad 1990 (symbol grounding). The maintenance
  logic of many-to-many symbol↔concept mappings follows WordNet (Miller
  1990; Fellbaum 1998), Word Sense Disambiguation (Navigli 2009), and
  SKOS (Miles \& Bechhofer 2009). Our S-layer operationalizes ambiguity
  as first-class and context-convergent, as argued in our
  symbol--concept mapping survey.
\item
  \textbf{Syntax as a definable object and grammar engineering.}
  W-grammars (van Wijngaarden 1965; ALGOL 68 1975); parsing expression
  grammars (Ford 2004); invertible syntax descriptions unifying parsing
  and pretty-printing (Rendel \& Ostermann 2010); concrete syntax as
  typed code with explicit decoders (our survey of
  syntax-as-definable-object). This motivates the native-grammar design:
  grammar itself is tokens.
\item
  \textbf{Variable binding and scope.} de Bruijn 1972 (nameless
  dummies); higher-order abstract syntax (Pfenning \& Elliott 1988;
  Harper--Honsell--Plotkin 1993 LF; McDowell \& Miller 2002); nominal
  logic (Gabbay \& Pitts 2002); locally nameless (Charguéraud 2012).
  Binding arity and scope verification (our G-layer) follow these;
  position/scope coverage is an independent degree of freedom.
\item
  \textbf{Systematic generalization.} Lake et al.~2017; Keysers et
  al.~2020 (COGS); SCAN. These document the phenomenon; this work
  provides the genesis (why and when generalization holds) as formalized
  induction.
\item
  \textbf{Neural-symbolic arithmetic modules.} NALU/NAU (fixed operator
  sets); this work uses a unified parameter set with zero-sample
  semantic closure for new operators.
\item
  \textbf{Length generalization.} Position-encoding and
  recurrent-elaboration engineering under fixed representations; this
  work removes length-dependence at the representation layer.
\item
  \textbf{Representation learning.} Disentangled/compositional
  structured representations study what representations look like; this
  work studies how representation determines learnability, treating
  representation design as a contribution.
\item
  \textbf{Program-to-neural compilation.} Learning to Compile Programs
  to Neural Networks (ICML 2024); inference compilation; shortcut
  distillation. Differences: same-model construction-channel
  self-compilation, no semantic-closure expansion, verifier + fallback.
\item
  \textbf{Dual-process cognition.} Kahneman (System 1/2) grounds G5:
  intuition as the reasoning fast path.
\item
  \textbf{Philosophy of mathematics --- three-channel equivalence.} The
  mathematical structures are all prior art (de Bruijn 1972; combinatory
  logic; Bourbaki structuralism; category theory; Mizar/Lean/Isabelle)
  and are explicitly not claimed. The claim is only three-channel
  equivalence --- construction/intuition/formal rewriting mutually
  compilable within one system --- the neural implementation of
  Computational Trinitarianism (Curry--Howard: programs = proofs =
  types), equivalence between execution channels rather than between
  logics. Stance: anti-Platonist --- intuition is compiled construction,
  not innate.
\item
  \textbf{Anti-nativism.} Fodor (language of thought) and Chomsky
  (universal grammar) are the opposing positions; empirical
  counterevidence (definitionally-impoverished tokens do not generalize)
  is provided.
\end{enumerate}

\subsection{10 Limitations}\label{limitations}

\begin{enumerate}
\def\labelenumi{\arabic{enumi}.}
\tightlist
\item
  \textbf{Serialization (positional bias).} Position coverage is an
  independent degree of freedom: a syntactic slot never covered in
  training does not generalize even with complete definitions.
  Positional dependence is not fully removable in autoregressive
  generation.
\item
  \textbf{Finite resources.} Weight/sample/epoch bounds exist;
  convergence is config-dependent (3 simple, 8 full --- exp53);
  hyper-operator results can be astronomically large, so judgment
  negatives use small domains with overflow guards.
\item
  \textbf{Operator-order coverage is incomplete.} Hyper-operator levels
  (power/root/tetration) beyond the trained level are not fully
  evaluated: radix-base generation overflows for astronomical results,
  so only addition↔multiplication level extrapolation is measured. This
  is a coverage limitation, not a negative result. Likewise, symbol
  permutation (exp41) permuted digit concepts only; full operator-symbol
  permutation is pending.
\item
  \textbf{Supported calculus ≠ unrestricted mathematics.} Only
  structures with formalized definitions (within the five definition
  forms) are covered; no claim of general mathematics or autonomous
  concept discovery; new concepts still require manual definition
  registration.
\item
  \textbf{Value space/representation.} Currently boolean judgments and
  counting structures; the full value space (fractions/complex) is
  partially connected; the numeral-value side is not fully open.
\item
  \textbf{Intuition is statistical, not guaranteed.} The intuition
  channel performs structural location statistically; semantic
  equivalence requires a verifier/fallback (construction channel) as
  backing.
\end{enumerate}

\subsection{11 Conclusion}\label{conclusion}

\begin{itemize}
\tightlist
\item
  Generalization is the formalization of induction (§3): definition
  completeness (G1) plus relational sufficiency (G2) yields systematic
  generalization in a minimal model --- engineered, reproducible. Its
  per-concept learning curves are bimodal (compiled vs not), making
  failures attributable to definitions and coverage, not to training
  volume.
\item
  Representation is the contribution: factored degrees of freedom turn
  generalization into a low-complexity learning problem; generalization
  is written into the geometry of the representation, not left to the
  model to discover. The de-factoring control (EXP-C1b) confirms the
  causal role: hiding the iteration structure collapses exactly those
  operators whose definitions depend on it (root, tetration), while
  reducible ones (power) survive.
\item
  Weight compilation: training compiles construction into intuition;
  supervision excludes traces, so structure is learned. Intuition is the
  reasoning fast path (G5) --- hard to obtain, easy to generalize by
  skipping steps, with zero decay across width, radix, and
  anonymization.
\item
  Three-channel equivalence (R): construction/intuition/formal rewriting
  are mutually compilable in neural implementation (exp50: 2000-digit
  carry chain, all three channels 1.000) --- the engineering unification
  of a century-old split in the foundations of mathematics.
\item
  Future: larger domains; autonomous concept formation (inducing new
  definitions from data); cross-modality (extensions of symbol
  permutation).
\end{itemize}

\subsection{Appendix}\label{appendix}

\begin{itemize}
\tightlist
\item
  A. Full experiment tables: config/seed/judgment per-token
  shortfalls/conclusions (docs/paper\_data/results.csv).
\item
  B. Grammar: complete gtoken/ptoken listing (data-driven source);
  presentation examples (add\_infix, and\_infix, not\_prefix,
  succ\_prefix).
\item
  C. More examples: training sample sequences (numeral\_split, law,
  three-channel inputs, judgment sequences).
\item
  D. Training configs: complete JSON (exp02 suite: synth.samples
  composition + judge + train).
\item
  E. Lean proofs: definition-chain completeness (no dangling references,
  cycle certificates) and the compilation correctness of three-channel
  equivalence --- discipline verification, not a new mathematical claim.
\end{itemize}

\subsection{Cross-References within the Unified
Framework}\label{cross-references-within-the-unified-framework}

This paper is one leg of the \emph{Unified Framework of Representation,
Logic, and Intuition} (manifesto:
\texttt{Unified\_Framework\_Representation\_Logic\_Intuition.md}). The
three legs co-adapt:

\begin{longtable}[]{@{}
  >{\raggedright\arraybackslash}p{(\linewidth - 4\tabcolsep) * \real{0.3333}}
  >{\raggedright\arraybackslash}p{(\linewidth - 4\tabcolsep) * \real{0.3333}}
  >{\raggedright\arraybackslash}p{(\linewidth - 4\tabcolsep) * \real{0.3333}}@{}}
\toprule\noalign{}
\begin{minipage}[b]{\linewidth}\raggedright
Leg
\end{minipage} & \begin{minipage}[b]{\linewidth}\raggedright
Paper
\end{minipage} & \begin{minipage}[b]{\linewidth}\raggedright
Claim
\end{minipage} \\
\midrule\noalign{}
\endhead
\bottomrule\noalign{}
\endlastfoot
form → construction & P1 (this paper) & correct syntax → extreme
training speed; correct construction → extreme generalization \\
construction → intuition & P0 (three-channel equivalence) & construction
= formal = intuition; intuition is compiled construction \\
construction → intuition (method) & P2 (neural macro compilation) &
compile the slow construction path into the fast reasoning path \\
\end{longtable}

Cyclic closure: form directs construction (P1), construction earns
intuition (P0/P2), intuition searches form (the fast path extrapolates
and re-expresses in the unified form). The unified conclusion --- form
directs logical construction, construction earns flashing intuition,
intuition exhaustively searches forms --- is stated in the manifesto and
in each paper's Unified Conclusion.

\subsection{References}\label{references}

{[}Bibliography in progress; citations sourced from the project's
cross-disciplinary surveys (docs/语法调研/). Verified full records to be
completed.{]}

\begin{itemize}
\tightlist
\item
  Aczel, P. 1977. An Introduction to Inductive Definitions. Handbook of
  Mathematical Logic.
\item
  Adámek, J., Herrlich, H., Strecker, G. 2004. Abstract and Concrete
  Categories.
\item
  Barendregt, H. 1984. The Lambda Calculus: Its Syntax and Semantics.
\item
  Belnap, N. 1993. On Rigorous Definitions.
\item
  Charguéraud, A. 2012. The Locally Nameless Representation. JAR.
\item
  de Bruijn, N. G. 1972. Lambda Calculus Notation with Nameless Dummies.
\item
  Fiore, M., Plotkin, G., Turi, D. 1999. Abstract Syntax and Variable
  Binding. LICS.
\item
  Flach, J. M., Moreira, Á. F., Lamb, L. C. 2023. Towards a Neural
  Lambda Calculus: Neurosymbolic AI Applied to the Foundations of
  Functional Programming. arXiv:2304.09276.
\item
  Fodor, J. 1975. The Language of Thought.
\item
  Ford, B. 2004. Parsing Expression Grammars. POPL.
\item
  Frege, G. 1892. Über Sinn und Bedeutung.
\item
  Gabbay, M. J., Pitts, A. M. 2002. A New Approach to Abstract Syntax
  with Variable Binding. FAC.
\item
  Goguen, J., Thatcher, J., Wagner, E., Wright, J. 1977. Initial Algebra
  Semantics and Continuous Algebras. JACM.
\item
  Gupta, A. 2023. Definitions. Stanford Encyclopedia of Philosophy.
\item
  Gupta, A., Belnap, N. 1993. The Revision Theory of Truth.
\item
  Harper, R., Honsell, F., Plotkin, G. 1993. A Framework for Defining
  Logics. JACM.
\item
  Harnad, S. 1990. The Symbol Grounding Problem. Physica D.
\item
  Kahneman, D. 2011. Thinking, Fast and Slow.
\item
  Keysers, D. et al.~2020. Measuring Compositional Generalization: A
  Comprehensive Method. ICLR.
\item
  Lake, B., Ullman, T., Tenenbaum, J., Gershman, S. 2017. Building
  Machines That Learn and Think Like People. BBS.
\item
  Lambek, J. 1968. A Fixpoint Theorem for Complete Categories. Math. Z.
\item
  Mac Lane, S. 1971. Categories for the Working Mathematician.
\item
  Martin-Löf, P. 1984. Intuitionistic Type Theory. Bibliopolis.
\item
  McDowell, R., Miller, D. 2002. Reasoning with Higher-Order Abstract
  Syntax in a Logical Framework. ACM TOCL.
\item
  Miles, A., Bechhofer, S. 2009. SKOS Reference. W3C.
\item
  Miller, G. et al.~1990. Introduction to WordNet.
\item
  Navigli, R. 2009. Word Sense Disambiguation: A Survey. ACM CSUR.
\item
  Pfenning, F., Elliott, C. 1988. Higher-Order Abstract Syntax. PLDI.
\item
  Rendel, T., Ostermann, K. 2010. Invertible Syntax Descriptions.
  Haskell Symposium.
\item
  Saussure, F. de. 1916. Cours de linguistique générale.
\item
  Suppes, P. 1957. Introduction to Logic.
\item
  van Wijngaarden, A. 1965. Orthogonal Design and Description of a
  Formal Language.
\item
  {[}Additional survey-derived records: Burris \& Sankappanavar 1981
  (Universal Algebra); Kripke 1963 (Modal Logic); Hamblin 1973
  (Questions); Groenendijk \& Stokhof 1984; Kamp 1981 (DRT); Ogden \&
  Richards 1923 (The Meaning of Meaning); Cruse 2000; Kilgarriff 1997;
  Cajori 1928 (A History of Mathematical Notations); Putnam 1975 (The
  Meaning of ``Meaning''); Rosch 1973 (Natural Categories); Enderton
  2001; Hodges 1993 (Model Theory); Gentzen 1934/35; Prawitz 1965;
  Russell 1905 (On Denoting); Li \& Vitányi 2008 (Kolmogorov
  Complexity); Fiore \& Szamozvancev 2022; Crary \& Harper 2006.{]}
\end{itemize}

\begin{center}\rule{0.5\linewidth}{0.5pt}\end{center}

\subsection{Unified Conclusion (added
2026-08-11)}\label{unified-conclusion-added-2026-08-11}

Once, under a unified formal attention-syntactic structure, a neural
network ultimately, inevitably, and necessarily clearly, intuitively,
and interpretably observes the boundary between intuition and
construction, we can with full confidence use the unified
attention-formal language to direct logical construction, use precise
synthetic-CoT logical construction to earn \emph{flashing intuition},
and use distilled intuitive context to exhaustively search attention
forms. Form, construction, and intuition co-adapt within the unified
framework, forming a feedback iteration pipeline of teaching, training,
and reasoning in which they are no longer distinct --- perhaps one of
the faster shortcuts among the many paths toward ASI.

\end{document}
